\documentclass[10pt,a4paper]{amsart}
\usepackage{amsmath,amssymb,amsthm,graphicx}
% Retained minimal preamble for the Stacks Project chapter flat.tex (More on Flatness).
% Verbatim excerpt of the chapter preamble: Xy-pic setup (native preamble lines 14-18) and the native
% theorem declarations with the mathematical macros (native preamble lines 154-217). Copied mechanically;
% nothing added, removed or redefined. The wrapper class and the amsmath/amssymb/amsthm/graphicx base
% packages are supplied by the assembler, so the class/externaldocument part of the native preamble is omitted.
% For commutative diagrams we use Xy-pic
\usepackage[all]{xy}

% We use 2cell for 2-commutative diagrams.
\xyoption{2cell}
% Theorem environments.
%
\theoremstyle{plain}
\newtheorem{theorem}[subsection]{Theorem}
\newtheorem{proposition}[subsection]{Proposition}
\newtheorem{lemma}[subsection]{Lemma}

\theoremstyle{definition}
\newtheorem{definition}[subsection]{Definition}
\newtheorem{example}[subsection]{Example}
\newtheorem{exercise}[subsection]{Exercise}
\newtheorem{situation}[subsection]{Situation}

\theoremstyle{remark}
\newtheorem{remark}[subsection]{Remark}
\newtheorem{remarks}[subsection]{Remarks}

\numberwithin{equation}{subsection}

% Macros
%
\def\lim{\mathop{\mathrm{lim}}\nolimits}
\def\colim{\mathop{\mathrm{colim}}\nolimits}
\def\Spec{\mathop{\mathrm{Spec}}}
\def\Hom{\mathop{\mathrm{Hom}}\nolimits}
\def\Ext{\mathop{\mathrm{Ext}}\nolimits}
\def\SheafHom{\mathop{\mathcal{H}\!\mathit{om}}\nolimits}
\def\SheafExt{\mathop{\mathcal{E}\!\mathit{xt}}\nolimits}
\def\Sch{\mathit{Sch}}
\def\Mor{\mathop{\mathrm{Mor}}\nolimits}
\def\Ob{\mathop{\mathrm{Ob}}\nolimits}
\def\Sh{\mathop{\mathit{Sh}}\nolimits}
\def\NL{\mathop{N\!L}\nolimits}
\def\CH{\mathop{\mathrm{CH}}\nolimits}
\def\proetale{{pro\text{-}\acute{e}tale}}
\def\etale{{\acute{e}tale}}
\def\QCoh{\mathit{QCoh}}
\def\Ker{\mathop{\mathrm{Ker}}}
\def\Im{\mathop{\mathrm{Im}}}
\def\Coker{\mathop{\mathrm{Coker}}}
\def\Coim{\mathop{\mathrm{Coim}}}

% Boxtimes
%
\DeclareMathSymbol{\boxtimes}{\mathbin}{AMSa}{"02}

%
% Macros for moduli stacks/spaces
%
\def\QCohstack{\mathcal{QC}\!\mathit{oh}}
\def\Cohstack{\mathcal{C}\!\mathit{oh}}
\def\Spacesstack{\mathcal{S}\!\mathit{paces}}
\def\Quotfunctor{\mathrm{Quot}}
\def\Hilbfunctor{\mathrm{Hilb}}
\def\Curvesstack{\mathcal{C}\!\mathit{urves}}
\def\Polarizedstack{\mathcal{P}\!\mathit{olarized}}
\def\Complexesstack{\mathcal{C}\!\mathit{omplexes}}
% \Pic is the operator that assigns to X its picard group, usage \Pic(X)
% \Picardstack_{X/B} denotes the Picard stack of X over B
% \Picardfunctor_{X/B} denotes the Picard functor of X over B
\def\Pic{\mathop{\mathrm{Pic}}\nolimits}
\def\Picardstack{\mathcal{P}\!\mathit{ic}}
\def\Picardfunctor{\mathrm{Pic}}
\def\Deformationcategory{\mathcal{D}\!\mathit{ef}}

% --- appended 2026-10-12 (Chief): the project own definitions of the reference/slogan/history
% display environments, copied mechanically from the shared Stacks preamble lines 8-12
% (\usepackage{verbatim} plus the three \newenvironment lines, which the chapter preamble keeps
% verbatim from the project preamble). Nothing added, removed or redefined: these environments
% carry no mathematical content and the project itself typesets them as comments. Without these
% three lines a counted statement whose own text uses them fails with Environment slogan undefined.
\usepackage{verbatim}
\newenvironment{reference}{\comment}{\endcomment}
\newenvironment{slogan}{\comment}{\endcomment}
\newenvironment{history}{\comment}{\endcomment}
\usepackage{cleveref}
\crefname{originalreference}{original reference}{original references}
\Crefname{originalreference}{Original reference}{Original references}
\makeatletter
\newlabel{situation-mod-injective}{{0AT0}{source}{}{sourcepointer-1}{}}
\newlabel{situation-mod-injective@cref}{{[originalreference][0][]0AT0}{[1][0][]source}{}{}{}}
\makeatother
\crefname{originallemma}{original lemma}{original lemmas}
\Crefname{originallemma}{Original lemma}{Original lemmas}
\makeatletter
\newlabel{more-algebra-lemma-valuation-ring-torsion-free-flat}{{0539}{source}{}{sourcepointer-2}{}}
\newlabel{more-algebra-lemma-valuation-ring-torsion-free-flat@cref}{{[originallemma][0][]0539}{[1][0][]source}{}{}{}}
\makeatother
\crefname{originallemma}{original lemma}{original lemmas}
\Crefname{originallemma}{Original lemma}{Original lemmas}
\makeatletter
\newlabel{more-algebra-lemma-flat-finite-type-valuation-ring-finite-presentation}{{053E}{source}{}{sourcepointer-3}{}}
\newlabel{more-algebra-lemma-flat-finite-type-valuation-ring-finite-presentation@cref}{{[originallemma][0][]053E}{[1][0][]source}{}{}{}}
\makeatother
\crefname{originallemma}{original lemma}{original lemmas}
\Crefname{originallemma}{Original lemma}{Original lemmas}
\makeatletter
\newlabel{algebra-lemma-extension}{{0519}{source}{}{sourcepointer-4}{}}
\newlabel{algebra-lemma-extension@cref}{{[originallemma][0][]0519}{[1][0][]source}{}{}{}}
\makeatother
\crefname{originallemma}{original lemma}{original lemmas}
\Crefname{originallemma}{Original lemma}{Original lemmas}
\makeatletter
\newlabel{algebra-lemma-flat-tor-zero}{{00HL}{source}{}{sourcepointer-5}{}}
\newlabel{algebra-lemma-flat-tor-zero@cref}{{[originallemma][0][]00HL}{[1][0][]source}{}{}{}}
\makeatother
\crefname{originallemma}{original lemma}{original lemmas}
\Crefname{originallemma}{Original lemma}{Original lemmas}
\makeatletter
\newlabel{algebra-lemma-NAK}{{00DV}{source}{}{sourcepointer-6}{}}
\newlabel{algebra-lemma-NAK@cref}{{[originallemma][0][]00DV}{[1][0][]source}{}{}{}}
\makeatother
\crefname{originaltheorem}{original theorem}{original theorems}
\Crefname{originaltheorem}{Original theorem}{Original theorems}
\makeatletter
\newlabel{algebra-theorem-universally-exact-criteria}{{058K}{source}{}{sourcepointer-7}{}}
\newlabel{algebra-theorem-universally-exact-criteria@cref}{{[originaltheorem][0][]058K}{[1][0][]source}{}{}{}}
\makeatother
\crefname{originallemma}{original lemma}{original lemmas}
\Crefname{originallemma}{Original lemma}{Original lemmas}
\makeatletter
\newlabel{more-algebra-lemma-generalized-valuation-ring-modules}{{0ASP}{source}{}{sourcepointer-8}{}}
\newlabel{more-algebra-lemma-generalized-valuation-ring-modules@cref}{{[originallemma][0][]0ASP}{[1][0][]source}{}{}{}}
\makeatother
\crefname{originallemma}{original lemma}{original lemmas}
\Crefname{originallemma}{Original lemma}{Original lemmas}
\makeatletter
\newlabel{lemma-flat-finite-type-local-valuation-ring-has-content}{{0ASX}{source}{}{sourcepointer-9}{}}
\newlabel{lemma-flat-finite-type-local-valuation-ring-has-content@cref}{{[originallemma][0][]0ASX}{[1][0][]source}{}{}{}}
\makeatother
\crefname{originallemma}{original lemma}{original lemmas}
\Crefname{originallemma}{Original lemma}{Original lemmas}
\makeatletter
\newlabel{more-algebra-lemma-content-finitely-generated}{{0ASB}{source}{}{sourcepointer-10}{}}
\newlabel{more-algebra-lemma-content-finitely-generated@cref}{{[originallemma][0][]0ASB}{[1][0][]source}{}{}{}}
\makeatother
\crefname{originallemma}{original lemma}{original lemmas}
\Crefname{originallemma}{Original lemma}{Original lemmas}
\makeatletter
\newlabel{algebra-lemma-characterize-valuation-ring}{{090Q}{source}{}{sourcepointer-11}{}}
\newlabel{algebra-lemma-characterize-valuation-ring@cref}{{[originallemma][0][]090Q}{[1][0][]source}{}{}{}}
\makeatother
\crefname{originallemma}{original lemma}{original lemmas}
\Crefname{originallemma}{Original lemma}{Original lemmas}
\makeatletter
\newlabel{more-algebra-lemma-equal-content}{{0ASC}{source}{}{sourcepointer-12}{}}
\newlabel{more-algebra-lemma-equal-content@cref}{{[originallemma][0][]0ASC}{[1][0][]source}{}{}{}}
\makeatother
\crefname{originaldefinition}{original definition}{original definitions}
\Crefname{originaldefinition}{Original definition}{Original definitions}
\makeatletter
\newlabel{more-algebra-definition-content-ideal}{{0ASA}{source}{}{sourcepointer-13}{}}
\newlabel{more-algebra-definition-content-ideal@cref}{{[originaldefinition][0][]0ASA}{[1][0][]source}{}{}{}}
\makeatother
\title{The injectivity lemma holds over a valuation ring: the module is flat and finitely presented and the map is universally injective}
\author{The Stacks Project Authors}
\date{stacks-project:flat.tex:local-snapshot-20261011; GNU FDL 1.2 or later}
\begin{document}
\maketitle
\noindent\textit{Editorial note.} All delivered mathematical text is the author's own, copied from the pinned Stacks Project chapter file flat.tex; no statement or proof step was rewritten, reordered or removed.\par\medskip
\noindent\textit{Editorial note.} The preamble is the chapter's own retained minimal preamble (sources/stacks-flat/retained-preamble.tex), a verbatim excerpt of the native Xy-pic setup and of the native theorem declarations and mathematical macros, followed by the project own definitions of the reference, slogan and history display environments (verbatim package plus three newenvironment lines) that the retained file carries at its end and that the project itself typesets as comments. No macro was invented or redefined, and the wrapper class and the amsmath/amssymb/amsthm/graphicx base packages are supplied by the assembler.\par\medskip
\noindent\textit{Editorial note.} The counted statement and proof contain no cross reference, no citation and no included file, so the delivered body is self contained.\par\medskip
\noindent\textit{Editorial note.} Canonical location of the counted proposition: Stacks tag 0AT5 (sources/stacks-flat/pinned-original-source.tex line 7158).\par\medskip
\noindent\textit{Allowed standard prerequisites (editorial).}\begin{itemize}
\item The situation and its three assertions. Same chapter, Situation labelled situation-mod-injective, invoked in the statement. Supplies the data A, B, N, M, u and the three properties to be proved.
\item Flatness equals torsion freeness over a valuation ring, and finite presentation of flat finite type modules. More on Algebra, Lemmas labelled more-algebra-lemma-valuation-ring-torsion-free-flat and more-algebra-lemma-flat-finite-type-valuation-ring-finite-presentation, cited at proof lines 7163-7172. Identifies flatness with torsion freeness and upgrades flatness of N/T to finite presentation over B.
\item The criterion for universal injectivity and the reduction of test modules. Algebra, Theorem labelled algebra-theorem-universally-exact-criteria, cited at proof lines 7186-7188; More on Algebra, Lemma labelled more-algebra-lemma-generalized-valuation-ring-modules, cited at proof lines 7188-7191. Reduces universal injectivity to injectivity after tensoring with A/(a).
\item Content ideals over a valuation ring. Same chapter, Lemma labelled lemma-flat-finite-type-local-valuation-ring-has-content, cited at proof lines 7195-7197; More on Algebra, Lemmas labelled more-algebra-lemma-content-finitely-generated and more-algebra-lemma-equal-content and Definition labelled more-algebra-definition-content-ideal, cited at proof lines 7197-7203. Produces the principal content ideal of x, identifies it with the content ideal of u(x), and compares it with (a).
\end{itemize}\medskip
\ifdefined\hypertarget\hypertarget{sourcepointer-1}{}\fi
\noindent\textit{Source reference.} Original reference 0AT0: sources/stacks-flat/pinned-original-source.tex, source lines 7017--7017. Same-chapter pointer into the author environment situation at source line 7017; invoked at proof line None of the counted proof. It is retained only as a navigation pointer; every load-bearing step is inside the delivered span.\par\medskip
\ifdefined\hypertarget\hypertarget{sourcepointer-2}{}\fi
\noindent\textit{Source reference.} Original lemma 0539: sources/stacks-chow/tags.txt, source lines 6597--6597. Cross-chapter standard prerequisite located in the project tag list staged in the local snapshot; invoked at proof line 7166 of the counted proof. It is retained only as a navigation pointer; every load-bearing step is inside the delivered span.\par\medskip
\ifdefined\hypertarget\hypertarget{sourcepointer-3}{}\fi
\noindent\textit{Source reference.} Original lemma 053E: sources/stacks-chow/tags.txt, source lines 6602--6602. Cross-chapter standard prerequisite located in the project tag list staged in the local snapshot; invoked at proof line 7170 of the counted proof. It is retained only as a navigation pointer; every load-bearing step is inside the delivered span.\par\medskip
\ifdefined\hypertarget\hypertarget{sourcepointer-4}{}\fi
\noindent\textit{Source reference.} Original lemma 0519: sources/stacks-chow/tags.txt, source lines 6527--6527. Cross-chapter standard prerequisite located in the project tag list staged in the local snapshot; invoked at proof line 7172 of the counted proof. It is retained only as a navigation pointer; every load-bearing step is inside the delivered span.\par\medskip
\ifdefined\hypertarget\hypertarget{sourcepointer-5}{}\fi
\noindent\textit{Source reference.} Original lemma 00HL: sources/stacks-chow/tags.txt, source lines 647--647. Cross-chapter standard prerequisite located in the project tag list staged in the local snapshot; invoked at proof line 7175 of the counted proof. It is retained only as a navigation pointer; every load-bearing step is inside the delivered span.\par\medskip
\ifdefined\hypertarget\hypertarget{sourcepointer-6}{}\fi
\noindent\textit{Source reference.} Original lemma 00DV: sources/stacks-chow/tags.txt, source lines 512--512. Cross-chapter standard prerequisite located in the project tag list staged in the local snapshot; invoked at proof line 7178 of the counted proof. It is retained only as a navigation pointer; every load-bearing step is inside the delivered span.\par\medskip
\ifdefined\hypertarget\hypertarget{sourcepointer-7}{}\fi
\noindent\textit{Source reference.} Original theorem 058K: sources/stacks-chow/tags.txt, source lines 6783--6783. Cross-chapter standard prerequisite located in the project tag list staged in the local snapshot; invoked at proof line 7184 of the counted proof. It is retained only as a navigation pointer; every load-bearing step is inside the delivered span.\par\medskip
\ifdefined\hypertarget\hypertarget{sourcepointer-8}{}\fi
\noindent\textit{Source reference.} Original lemma 0ASP: sources/stacks-chow/tags.txt, source lines 13587--13587. Cross-chapter standard prerequisite located in the project tag list staged in the local snapshot; invoked at proof line 7188 of the counted proof. It is retained only as a navigation pointer; every load-bearing step is inside the delivered span.\par\medskip
\ifdefined\hypertarget\hypertarget{sourcepointer-9}{}\fi
\noindent\textit{Source reference.} Original lemma 0ASX: sources/stacks-flat/pinned-original-source.tex, source lines 5801--5801. Same-chapter pointer into the author environment lemma at source line 5801; invoked at proof line 7192 of the counted proof. It is retained only as a navigation pointer; every load-bearing step is inside the delivered span.\par\medskip
\ifdefined\hypertarget\hypertarget{sourcepointer-10}{}\fi
\noindent\textit{Source reference.} Original lemma 0ASB: sources/stacks-chow/tags.txt, source lines 13574--13574. Cross-chapter standard prerequisite located in the project tag list staged in the local snapshot; invoked at proof line 7195 of the counted proof. It is retained only as a navigation pointer; every load-bearing step is inside the delivered span.\par\medskip
\ifdefined\hypertarget\hypertarget{sourcepointer-11}{}\fi
\noindent\textit{Source reference.} Original lemma 090Q: sources/stacks-chow/tags.txt, source lines 11418--11418. Cross-chapter standard prerequisite located in the project tag list staged in the local snapshot; invoked at proof line 7197 of the counted proof. It is retained only as a navigation pointer; every load-bearing step is inside the delivered span.\par\medskip
\ifdefined\hypertarget\hypertarget{sourcepointer-12}{}\fi
\noindent\textit{Source reference.} Original lemma 0ASC: sources/stacks-chow/tags.txt, source lines 13575--13575. Cross-chapter standard prerequisite located in the project tag list staged in the local snapshot; invoked at proof line 7198 of the counted proof. It is retained only as a navigation pointer; every load-bearing step is inside the delivered span.\par\medskip
\ifdefined\hypertarget\hypertarget{sourcepointer-13}{}\fi
\noindent\textit{Source reference.} Original definition 0ASA: sources/stacks-chow/tags.txt, source lines 13573--13573. Cross-chapter standard prerequisite located in the project tag list staged in the local snapshot; invoked at proof line 7201 of the counted proof. It is retained only as a navigation pointer; every load-bearing step is inside the delivered span.\par\medskip
\begin{lemma}
\label{lemma-valuation-ring-finite-type-injective-into-flat-mod-m}
If in Situation \ref{situation-mod-injective} the ring $A$ is a
valuation ring then the lemma holds.
\end{lemma}

\begin{proof}
Recall that an $A$-module is flat if and only if it is torsion free, see
More on Algebra, Lemma
\ref{more-algebra-lemma-valuation-ring-torsion-free-flat}.
Let $T \subset N$ be the $A$-torsion. Then $u(T) = 0$ and
$N/T$ is $A$-flat. Hence $N/T$ is finitely presented over $B$, see
More on Algebra, Lemma
\ref{more-algebra-lemma-flat-finite-type-valuation-ring-finite-presentation}.
Thus $T$ is a finite $B$-module, see
Algebra, Lemma \ref{algebra-lemma-extension}.
Since $N/T$ is $A$-flat we see that
$T/\mathfrak m_A T \subset N/\mathfrak m_A N$, see
Algebra, Lemma \ref{algebra-lemma-flat-tor-zero}.
As $\overline{u}$ is injective but $u(T) = 0$, we conclude that
$T/\mathfrak m_A T = 0$. Hence $T = 0$ by Nakayama's lemma, see
Algebra, Lemma \ref{algebra-lemma-NAK}. At this point we have
proved two out of the three assertions ($N$ is $A$-flat and
of finite presentation over $B$) and what is left is to show that
$u$ is universally injective.

\medskip\noindent
By Algebra, Theorem \ref{algebra-theorem-universally-exact-criteria}
it suffices to show that $N \otimes_A Q \to M \otimes_A Q$ is injective
for every finitely presented $A$-module $Q$. By
More on Algebra, Lemma
\ref{more-algebra-lemma-generalized-valuation-ring-modules}
we may assume $Q = A/(a)$ with $a \in \mathfrak m_A$ nonzero.
Thus it suffices to show that $N/aN \to M/aM$ is injective.
Let $x \in N$ with $u(x) \in aM$. By
Lemma \ref{lemma-flat-finite-type-local-valuation-ring-has-content}
we know that $x$ has a content ideal $I \subset A$. Since
$I$ is finitely generated
(More on Algebra, Lemma \ref{more-algebra-lemma-content-finitely-generated})
and $A$ is a valuation ring, we have $I = (b)$ for some $b$
(by Algebra, Lemma \ref{algebra-lemma-characterize-valuation-ring}).
By More on Algebra, Lemma \ref{more-algebra-lemma-equal-content}
the element $u(x)$ has content ideal $I$ as well.
Since $u(x) \in aM$ we see that $(b) \subset (a)$
by More on Algebra, Definition \ref{more-algebra-definition-content-ideal}.
Since $x \in bN$ we conclude $x \in aN$ as desired.
\end{proof}

\end{document}
